\HMTNarrativeHeading{What a transformation preserves}

Operator theory gives a calculable form to the conservation of the state
and its registers. A partial reading obtains one component; the
complement collects the contribution that allows the complete balance
to be restored. The sum of the squared norms of the component read and
its complement reproduces the squared norm of the initial state. The
recovery formulas specify which data and which operators make it
possible to invert the reading.

As transport is iterated, each stage incorporates a complement and
prolongs a terminal component. The finite balance explicitly preserves
this terminal component. The passage to arbitrary depth studies
convergence and determines when the archive accounts for the entire
norm. This order makes it possible to interpret memory as an operation
of conservation and recovery, with verifiable conditions.

The exterior extension shows how areas and volumes of the linear
realization are transported. Mixed contributions participate in the
exterior products and their balances. The change of representation
preserves the relations measured by these objects. Recovery maintains
the link with the orientation and forms used by each proof.

Monodromy and loops between geometrical regimes add a definite
dynamics. The article composes elliptic, parabolic, and hyperbolic
transformations, preserves their memory, and calculates their response.
The symmetries of the transport action produce the variational charges
of that domain. The invariance of an observable, in turn, is formulated
through its relation to transport. These two forms of conservation are
linked while preserving their particular definitions.

The physical realizations subsequently recover the action sections,
area readers, and scales generated previously. The Gaussian example
makes it possible to calculate the purity and spectrum of a reduced
reading from a specified joint preparation. Its correlations preserve
the information of the complete state. This result shows an operational
consequence of the argument: transport determines both what an
observation receives and the information needed to understand its
relation to the whole.
